% LuaLaTeX-ja を用いた卒業論文テンプレート（内容更新版）
\documentclass[a4paper,11pt]{ltjsreport}

%======================= パッケージ =======================
\usepackage{luatexja}
\usepackage{amsmath,amssymb}
\usepackage{bm}
\usepackage{graphicx}
\usepackage{ascmac}
\usepackage[top=20mm,bottom=30mm,left=20mm,right=20mm]{geometry}
\usepackage{multirow}
\usepackage{url}
\usepackage{color}
\renewcommand{\bibname}{参考文献}

%======================= 本文開始 =========================
\begin{document}

%--------------------- 表紙 ---------------------
\begin{titlepage}
  \centering
  ~~\vspace{1cm}
 
  {\Huge 卒　業　研　究　論　文 \par}
  {\Huge （2025年度）\par}

  \vspace{8cm}

  \begin{tabular}{rl}
    {\LARGE 題　　目}　 & {\LARGE ：低コスト移動ロボットにおける} \\
                         & {\LARGE ~~VLM単体制御によるターゲット追従の} \\
                         & {\LARGE ~~性能評価} \\ \\ \\
    {\LARGE 指導教員}　 & {\LARGE ：保坂 忠明} \\　 \\ \\
    {\LARGE 学生番号}　 & {\LARGE ：1512221209} \\　 \\ \\
    {\LARGE 氏　　名}　 & {\LARGE ：小坂 尚璃}
  \end{tabular}

  \vfil
  \begin{flushright}
    明治大学理工学部 電気電子生命学科
  \end{flushright}
\end{titlepage}

\tableofcontents

%===========================================================
\chapter{序論}
\label{chap:intro}
\section{研究背景}
近年，画像と言語を同時に扱う大規模モデル（本稿では Vision-Language Model，VLM と呼ぶ）が普及し，ロボットの目標物認識や移動制御への応用が進んでいる。本研究では VLM を唯一の知覚・意思決定モジュールとして用いる構成を前提とし，外部地図や追加センサに依存しない低コスト移動ロボットの実現性を検討する。教育機関における実験機材や，施設巡回・簡易搬送といった応用分野では，低コストで事前の環境地図作成を必要としない即応性の高い移動ロボットが求められる。一方で，視覚ノイズや照度変化により進行方向が乱れやすく，推論遅延が大きい場合は過去フレーム由来の誤判断が累積する弱点が残る。低コスト移動ロボットではこのシンプルな構成が現実的であり，性能限界を把握することが実装指針として重要である。

\section{先行研究と課題}
SayCan は大規模言語モデルを用いた行動計画を提案し \cite{saycan}，PaLM-E はマルチモーダル化で計画精度を高めた \cite{palme}。RT-2 は画像と言語から直接行動を出力し \cite{rt2}，VL-Maps は意味情報付き地図で探索性能を向上させた \cite{vlmaps}。これらは地図や外部記憶を併用するため高性能だが，準備時間と計算資源を要し，Raspberry Pi のような低価格プラットフォームへの適用は容易でない。外部地図を用いず低コストで VLM を評価した報告は少なく，走行安定性に関わるボトルネックは十分に整理されていない。教育・実験向けに普及している小型ロボットでは，推論を外部 PC にオフロードする構成が現実的であり，その場合の遅延と視覚品質の影響を測る必要がある。

\section{研究目的と評価設定}
本研究の目的は，外部地図や追加センサに依存せず VLM 単体で走行を決定する低コスト移動ロボットの性能を定量化し，失敗要因を明確にすることである。具体的には，OSOYOO Raspberry Pi Car を用い，単一カメラ映像を PC 側で推論して「黄色いターゲットボックス（以下，ターゲット）に接近し，前方 20 cm 以内で停止する」タスクを実環境で評価する。性能指標は到達時間，成功率，コマンド適合率，反転回数，推論遅延とし，これらの測定を通じて VLM 単体制御の限界と改善優先度を整理する。

\section{論文構成}
本論文の構成を示す。第\ref{chap:method}章でハードウェアと制御条件を述べる。第\ref{chap:experiments}章で実験環境と評価指標を説明し，第\ref{chap:results}章で暫定結果を示す。第\ref{chap:discussion}章で考察を行い，第\ref{chap:conclusion}章で結論と今後の課題を述べる。付録では補助的な試行結果を提示する。

%===========================================================
\chapter{手法}
\label{chap:method}

\section{ハードウェア構成}
対象機体は OSOYOO Raspberry Pi Robot Car（2023版）である。Raspberry Pi 4B，広角 USB カメラ，モータドライバ（L298N），ステアリング用サーボから構成し，カメラ解像度は 640×480，フレームレートは 5 fps とする。走行は室内の平坦な床面で行う。

\section{通信とソフトウェア構成}
Raspberry Pi 上で `raspi\_agent.py` を 8080 番ポートで起動し，`GET /frame.jpg` と `/stream.mjpg` によるフレーム配信，`POST /command` による `FORWARD/LEFT/RIGHT/BACKWARD/STOP` コマンド受信を提供する。PC 側は同一ネットワークから HTTP でアクセスし，VLM の推論結果に基づきコマンドを送信する。

\section{使用モデルとプロンプト}
主モデルは GPT-5-mini-20250807（ローカル実行）である。単一フレームを入力し，「ターゲットに到達して停止せよ」と指示した上で，選択肢を `LEFT/RIGHT/FORWARD/FORWARD\_SLOW/STOP/BACKWARD` の一語に限定する。

\section{制御条件}
\begin{itemize}
  \item \textbf{条件A（ダイレクト制御）}: VLM が各ループで移動コマンドを直接選び，そのままロボットへ送信する。推論に 5--7 秒要するため，ループ間隔は 10 秒とし，速度は固定とする。
  \item \textbf{条件B（ハイブリッド制御）}: VLM は「ターゲット位置（LEFT/CENTER/RIGHT）と距離（FAR/MID/NEAR）」のみを JSON で返し，PC 側の規則でコマンドを決定する。推論遅延を考慮しループ間隔は 10 秒に揃える。
\end{itemize}
両条件とも VLM 出力と単純規則のみで挙動を決める。

\begin{table}[htbp]
  \centering
  \caption{ハイブリッド制御の規則}
  \label{tab:hybrid-rule}
  \begin{tabular}{|c|c|c|}
    \hline
    距離 & 位置 & 送信コマンド \\
    \hline
    FAR/MID & CENTER & FORWARD \\
    NEAR    & CENTER & FORWARD\_SLOW または STOP \\
    任意    & LEFT   & LEFT \\
    任意    & RIGHT  & RIGHT \\
    \hline
  \end{tabular}
\end{table}

\section{比較観点と改善優先度}
ダイレクト制御とハイブリッド制御を比較し，誤りが生じやすい工程を二段階に分けて評価する。まず，VLM が画像から位置と距離を正しく判断できるかをハイブリッド制御で測定し，誤認識が多い場合は入力前処理や照度制御を優先して改善する。次に，判断は正しいがコマンド生成で逆転や過大旋回が起きる場合は，ダイレクト制御のプロンプト設計やコマンド選択肢の絞り込みを優先する。この結果から，外部記憶の導入や予測モデルの追加が必要となる箇所（判断工程か指令工程か）を明確にし，拡張設計の順序を定める。

%===========================================================
\chapter{実験}
\label{chap:experiments}

\section{環境とコース}
室内直線コース（長さ 2.5 m）と L 字コース（1.5 m＋1.5 m）を作成し，照度は 400--500 lx に保つ。壁や障害物は固定し，背景テクスチャを一定にする。ターゲットは黄色いターゲットボックス（“TARGET” と印字）を終点に配置する。L 字では角手前 0.4--0.6 m に赤い補助ボックスを置き，英字で “TURN” と印字する（ラベルは「TURN」のみとし，VLM プロンプトでは「赤い箱に到達したら左折，その後黄色い“TARGET”を再探索」と指示する）。照度ロバスト性を確認する簡易環境変化として，時間が許す場合は 300 lx（暗め）と 700 lx（明るめ）を各 2 試行追加する。

\section{開始姿勢と初期視野}
走行の再現性を確保するため，スタート位置に 1 m 四方のテープ枠と進行方向基準線を貼付する。プレビュー画面でターゲットを水平方向は中央±10\%，垂直方向は中央よりやや上に収め，画面幅の 5--15\% の大きさに見える距離で開始する。露出が安定するまで 1--2 秒静止し，ログ開始後に走行を開始する。頑健性確認として，時間が許す場合に左右 20--25\% ずれた初期位置で各 2 試行のみ実施し，その他の設定は固定する。

\section{条件設定と試行計画}
ダイレクト制御とハイブリッド制御を，直線コースと L 字コースの双方で比較する。各条件の目的は以下の通りである。
\begin{itemize}
  \item 条件D-S（ダイレクト・直線）: 基本性能と照度影響のベースラインを取得する。
  \item 条件H-S（ハイブリッド・直線）: 距離推定付きの速度切替が到達時間と反転回数をどこまで安定化させるかを測る。
  \item 条件D-L（ダイレクト・L字）: サーボを用いずに，角手前 STOP と連続旋回のみで曲がり切れる下限性能を評価する。
  \item 条件H-L（ハイブリッド・L字）: 距離推定と減速を併用した場合の再捕捉成功率と横偏差を評価する。
\end{itemize}
時間が極端に限られる場合は，ダイレクト制御のみに統一し，直線と L 字を標準照度で各 6 試行，暗め照度で各 4 試行を行う簡略版を採用する。簡略版でも角手前 STOP と連続旋回（5--6 ループで約 90°）は必ず挿入し，コーナ再捕捉の成否を記録する。

\section{手順と試行数}
開始位置はターゲットから約 2.5 m 離し，進行方向を基準線に合わせる。各条件で直線・L 字とも 8--10 試行を基本とし，簡略版では 6 試行に短縮する。L 字では角手前で必ず STOP を 1 回挟み，その後に右旋回コマンドを連続 5--6 ループ送ってから再度直進させる。補助ボックス（“TURN”）を置く場合，STOP 位置は赤い箱の直前とし，赤い箱が視野中央に入ったことを確認した上で右旋回に移る。走行中はフレーム画像，VLM 出力，送信コマンド，タイムスタンプを保存し，別カメラで全体動画を記録する。試行には通し番号を付け，コース種別，照度，モデルバージョン，開始・終了時刻，初期ターゲット位置（中央／左寄せ／右寄せ）をメタ情報として残す。補助ボックス有無はメタ情報に必ず記録し，表\ref{tab:summary}では補助なし／ありを別行で示す。

\section{評価指標}
\begin{itemize}
  \item 到達時間：開始コマンド送信から停止判定までの時間。
  \item 成功率：ターゲット前 20 cm 以内で停止した試行の割合。
  \item コマンド適合率：全送信コマンド中，理想方向・速度に合致すると評価できる割合。
  \item コマンド反転回数：左右または前後の即時反転の回数。
  \item 推論遅延：VLM 応答時間の中央値と 95 パーセンタイル。
\end{itemize}

%===========================================================
\chapter{結果（暫定）}
\label{chap:results}
直近のログには試行ごとの到達可否，コマンド系列，推論遅延が保存されている。現在集計中であり，確定値を表\ref{tab:summary}に示す予定である。表では直線・L 字の双方について，条件D-S/H-S/D-L/H-L あるいは簡略版ダイレクトのみの平均到達時間，成功率，コマンド適合率，反転回数中央値，推論遅延中央値／95 パーセンタイルを掲載する。初期ターゲット位置のずらし試行は別列で報告し，照度変更（300 lx，700 lx）がある場合は併記する。コマンド時系列と軌跡オーバーレイを図\ref{fig:trajectory1}，図\ref{fig:trajectory2}として挿入する計画である。

% 予備の表・図枠（後で数値を埋める）
\begin{table}[htbp]
  \centering
  \caption{各条件の暫定集計（後日更新）}
  \label{tab:summary}
  \begin{tabular}{|c|c|c|c|c|c|}
    \hline
    条件 & コース & 到達時間[s] & 成功率 & 適合率 & 反転中央値 \\
    \hline
    ダイレクト & 直線 & -- & -- & -- & -- \\
    ダイレクト & L字 & -- & -- & -- & -- \\
    ハイブリッド & 直線 & -- & -- & -- & -- \\
    ハイブリッド & L字 & -- & -- & -- & -- \\
    \hline
  \end{tabular}
\end{table}

\begin{figure}[htbp]
  \centering
  \fbox{\rule{0.65\textwidth}{0.32\textwidth}}
  \caption{軌跡オーバーレイ（後日実データに置換）}
  \label{fig:trajectory1}
\end{figure}

\begin{figure}[htbp]
  \centering
  \fbox{\rule{0.65\textwidth}{0.32\textwidth}}
  \caption{コマンド系列と距離推定の時系列（後日実データに置換）}
  \label{fig:trajectory2}
\end{figure}

%===========================================================
\chapter{考察（予定）}
\label{chap:discussion}
予備観察では，ダイレクト制御は推論遅延と誤指示によりコマンド反転が増え，到達時間がばらつく傾向があった。ハイブリッド制御は距離推定の粗さが残るものの，左右誤指示の影響を抑え，成功率が高まると見込む。VLM 単体制御の主要課題は，(1) 応答遅延がループ間隔に近いこと，(2) 単一フレーム依存による瞬間的な誤認識，(3) 照度変動への感度である。これらに対し，フレームスタックによる時系列文脈の付与，簡易外部記憶による状態保持，軽量予測モデルによる補正が有効と考える。

%===========================================================
\chapter{結論}
\label{chap:conclusion}
本研究は，低コスト移動体で地図や複雑なセンサ統合を用いない構成において，VLM 単体制御が達成できる性能を評価する枠組みを示した。ダイレクト制御とハイブリッド制御を同条件で比較し，到達時間・成功率・コマンド適合率を中心指標，反転回数と遅延分布を補助指標として計測する計画を立てた。今後，集計したログを反映して結果と考察を更新し，外部記憶や予測モデルを組み込む際の具体的設計指針を示す。

%===========================================================
\chapter*{謝辞}
\addcontentsline{toc}{chapter}{謝辞}
本研究の遂行にあたり，実験環境整備に協力いただいた研究室メンバー各位に感謝する。計算資源と機材提供に謝意を表する。

%===========================================================
\bibliographystyle{junsrt}
\bibliography{mybibfile}

%===========================================================
\appendix
\chapter{SmolVLM 失敗試行の概要}
SmolVLM2-mlx を用いて同一タスクを試行したところ，推論時間は平均 0.8 s でループ間隔 0.5 s を上回った。その結果，過去フレームに基づく誤指示が頻発し，10 試行中の成功は 2 回に留まった。誤指示の内訳は左右の即時反転 12 件，停止不能 4 件であった。今後はモデルの軽量化またはフレーム間補間の導入を検討する。

\end{document}
